\chapter{Родитель relative-curvature junk quotient}
\label{ch:v10-h15-r2-relative-curvature-junk-quotient-parent}

Предыдущий аудит оставил единственный positive exact fork: выделить
relative endpoint curvature до взятия нормы. Теперь необходимо отделить
standard represented junk quotient от дополнительной graph-Hodge проекции.

Для минимальной узловой алгебры $\mathbb C^2$ и оператора ребра
\begin{equation}
 D=\begin{pmatrix}0&1\\1&0\end{pmatrix}
\end{equation}
точное универсальное исчисление даёт
\begin{equation}
 \rank\pi_D(\Omega^1)=2,
 \qquad \rank\pi_D(\Omega^2)=2,
 \qquad \rank J_D^2=0.
\end{equation}
Следовательно, rank quotient также равен двум. В нём независимо выживают
как суммарный класс $(1,1)$, так и относительный класс $(1,-1)$. Standard
junk не выбирает нужную разность и не объясняет удаление ordinary trace.

Двухузловой граф имеет incidence и Laplacian
\begin{equation}
 \partial=(1,-1),
 \qquad L_{P_2}=\partial^T\partial=
 \begin{pmatrix}1&-1\\-1&1\end{pmatrix}.
\end{equation}
Нормированный Hodge projector
\begin{equation}
 P_{\rm rel}=\frac12L_{P_2}
\end{equation}
идемпотентен, имеет rank $1$, убивает $(1,1)$ и сохраняет $(1,-1)$.
Его положительная форма имеет inertia $(0_-,1_0,1_+)$. Применение
relative readout к endpoint moment maps даёт
\begin{equation}
 (T,-B)\begin{pmatrix}1\\-1\end{pmatrix}=T+B=Q,
\end{equation}
то есть точный оператор rank $8$ без отрицательной полной trace metric.

Это положительный условный результат, но не результат junk calculus.
Оператор $P_{\rm rel}$ требует Hodge-структуры узлового графа, отсутствующей
в текущем represented action: inherited selector имеет rank $0/1$.
Кроме того, само odd auxiliary edge имеет inherited rank $0/8$. Поэтому
graph-Hodge construction пока задаёт новую связанную архитектуру, а не
вывод из существующего mapping cone.

\begin{theorem}[Junk no-go и условный graph-Hodge selector]
В минимальном двухузловом calculus degree-two junk равен нулю и не
различает суммарный и относительный endpoint-классы. Канонический projector
$L_{P_2}/2$ положителен и точно выбирает $Q$, но происходит из отдельной
graph-Hodge структуры, которая вместе с auxiliary edge ещё не унаследована.
\end{theorem}

\begin{nogocriterion}[Запрет переименования Hodge projection в junk]
Нельзя считать relative selector стандартным Connes quotient: прямой
расчёт даёт $J_D^2=0$. Для физического parent необходимо предъявить общий
носитель, на котором incidence, Hodge metric, odd auxiliary edge и
superconnection curvature являются частями одного заранее заданного
действия.
\end{nogocriterion}

\begin{protocol}\begin{enumerate}
\item узловая алгебра: \textbf{$\mathbb C^2$};
\item ranks $\Omega^1/\Omega^2/J^2$: \textbf{$2/2/0$};
\item standard quotient rank: \textbf{$2$};
\item junk выбирает relative class: \textbf{нет};
\item graph-Hodge projector rank: \textbf{$1$};
\item projector убивает sum и сохраняет difference: \textbf{точно};
\item relative readout: \textbf{$Q$, rank $8$};
\item inherited Hodge selector rank: \textbf{$0/1$};
\item inherited auxiliary edge rank: \textbf{$0/8$};
\item conditional/physical origin: \textbf{$4/6$ и $4/6$};
\item следующий гейт: \textbf{common parent relative Hodge auxiliary edge}.
\end{enumerate}\end{protocol}

Машинный сертификат сохранён в
\path{s2t_v10_particle_wrinkle_dislocation_callias_equal_charge_twist_m4_h15_r2_relative_curvature_junk_quotient_parent_origin_gate_results.json}.